\documentclass[11pt,a4paper]{article}

\usepackage[utf8]{inputenc}
\usepackage[T1]{fontenc}
\usepackage[french]{babel}
\usepackage[a4paper,margin=2.2cm]{geometry}
\usepackage{amsmath,amssymb,amsthm,mathtools}
\usepackage{lmodern}
\usepackage{enumitem}
\usepackage{hyperref}
\usepackage{fancyhdr}
\usepackage{xcolor}

\hypersetup{
    colorlinks=true,
    linkcolor=black,
    urlcolor=blue
}

\setlength{\parindent}{0pt}
\setlength{\parskip}{0.6em}

\pagestyle{fancy}
\fancyhf{}
\fancyhead[L]{Fiche d’exercices --- Nombres complexes}
\fancyhead[R]{Terminale }
\fancyfoot[C]{\thepage}

\newcommand{\R}{\mathbb{R}}
\newcommand{\C}{\mathbb{C}}

\begin{document}

\begin{center}
    {\LARGE \textbf{Fiche d’exercices --- Nombres complexes}}\\[0.3em]
    {\large De la terminale à la première année de prépa}
\end{center}

\hrule
\vspace{0.8em}

\textbf{Aide de lecture :}  
\textbf{APP} = application directe, \textbf{CALC} = calcul, \textbf{GEO} = géométrie, \textbf{PREP} = plus poussée.  
Les étoiles donnent un gradient de difficulté.

\tableofcontents

\newpage

\section*{I. Forme algébrique et calculs}
\addcontentsline{toc}{section}{I. Forme algébrique et calculs}

\textbf{Exercice 1 [APP $\star$]}\\
Pour chacun des nombres complexes suivants, donner sa partie réelle et sa partie imaginaire :
\[
z_1=3+2i,\qquad z_2=-5i,\qquad z_3=7,\qquad z_4=-2+4i.
\]

\textbf{Exercice 2 [APP $\star$]}\\
Parmi les nombres suivants, préciser ceux qui sont réels, ceux qui sont imaginaires purs :
\[
2,\qquad -3i,\qquad 4+0i,\qquad i\sqrt{5},\qquad 1-i.
\]

\textbf{Exercice 3 [CALC $\star\star$]}\\
Calculer sous forme algébrique :
\begin{enumerate}[label=\alph*)]
    \item \((2+3i)+(4-5i)\),
    \item \((1-2i)(3+i)\),
    \item \((2+i)^2\),
    \item \((1+i)^3\).
\end{enumerate}

\textbf{Exercice 4 [CALC $\star\star$]}\\
Mettre sous forme algébrique :
\begin{enumerate}[label=\alph*)]
    \item \(\dfrac{1}{2+i}\),
    \item \(\dfrac{3-2i}{1+i}\),
    \item \(\dfrac{2+3i}{1-2i}\).
\end{enumerate}

\textbf{Exercice 5 [CALC $\star\star$]}\\
Résoudre dans \(\C\) :
\[
z+2i=3-5i,
\qquad
(1-i)z=4+2i.
\]

\textbf{Exercice 6 [PREP $\star\star\star$]}\\
Soient \(z_1=a+ib\) et \(z_2=c+id\).  
Retrouver les formules de \(\Re(z_1z_2)\) et \(\Im(z_1z_2)\).

\section*{II. Conjugué, module et propriétés}
\addcontentsline{toc}{section}{II. Conjugué, module et propriétés}

\textbf{Exercice 7 [APP $\star$]}\\
Calculer le conjugué des nombres suivants :
\[
3+2i,\qquad -1+4i,\qquad 5,\qquad -7i.
\]

\textbf{Exercice 8 [APP $\star\star$]}\\
Montrer que :
\begin{enumerate}[label=\alph*)]
    \item \(z\in \R \Longleftrightarrow z=\overline z\),
    \item \(z\in i\R \Longleftrightarrow z=-\overline z\).
\end{enumerate}

\textbf{Exercice 9 [CALC $\star\star$]}\\
Calculer le module des nombres suivants :
\[
1+i,\qquad -3+4i,\qquad 2,\qquad -i\sqrt{3}.
\]

\textbf{Exercice 10 [CALC $\star\star$]}\\
Vérifier sur des exemples que
\[
z\overline z=|z|^2.
\]

\textbf{Exercice 11 [CALC $\star\star$]}\\
Montrer que pour tout \(z\in\C\),
\[
|\Re(z)|\le |z|,\qquad |\Im(z)|\le |z|.
\]

\textbf{Exercice 12 [PREP $\star\star\star$]}\\
Soit \(z\neq 0\). Montrer que
\[
\frac{1}{z}=\frac{\overline z}{|z|^2}.
\]

\textbf{Exercice 13 [PREP $\star\star\star$]}\\
Montrer que pour tous \(z,z'\in\C\),
\[
|zz'|=|z||z'|.
\]
Puis en déduire, pour \(z'\neq 0\),
\[
\left|\frac{z}{z'}\right|=\frac{|z|}{|z'|}.
\]

\section*{III. Interprétation géométrique}
\addcontentsline{toc}{section}{III. Interprétation géométrique}

\textbf{Exercice 14 [GEO $\star\star$]}\\
Donner l’affixe des points suivants :
\[
A(2,3),\qquad B(-1,4),\qquad C(0,-2).
\]

\textbf{Exercice 15 [GEO $\star\star$]}\\
Si \(A\) et \(B\) ont pour affixes \(z_A=1+2i\) et \(z_B=4-i\), déterminer l’affixe du vecteur \(\overrightarrow{AB}\).

\textbf{Exercice 16 [GEO $\star\star$]}\\
Déterminer l’affixe du milieu du segment \([AB]\) dans les cas suivants :
\begin{enumerate}[label=\alph*)]
    \item \(z_A=2+i,\ z_B=4-3i\),
    \item \(z_A=-1+5i,\ z_B=3+i\).
\end{enumerate}

\textbf{Exercice 17 [GEO $\star\star$]}\\
Si \(A\) et \(B\) ont pour affixes respectives \(z_A\) et \(z_B\), montrer que
\[
AB=|z_B-z_A|.
\]

\section*{IV. Inégalité triangulaire}
\addcontentsline{toc}{section}{IV. Inégalité triangulaire}

\textbf{Exercice 18 [APP $\star\star$]}\\
Calculer \(|z+z'|\), \(|z|+|z'|\) pour :
\[
z=1+i,\qquad z'=1-i.
\]
Comparer les deux quantités.

\textbf{Exercice 19 [DEM $\star\star\star$]}\\
Démontrer l’inégalité triangulaire :
\[
|z+z'|\le |z|+|z'|.
\]

\textbf{Exercice 20 [DEM $\star\star\star$]}\\
En déduire que pour tous \(z,z'\in\C\),
\[
\big||z|-|z'|\big|\le |z-z'|.
\]

\section*{V. Forme exponentielle et trigonométrique}
\addcontentsline{toc}{section}{V. Forme exponentielle et trigonométrique}

\textbf{Exercice 21 [APP $\star\star$]}\\
Écrire sous forme exponentielle :
\[
1,\qquad -1,\qquad i,\qquad -i.
\]

\textbf{Exercice 22 [APP $\star\star$]}\\
Déterminer le module et un argument des nombres suivants :
\[
z_1=1+i,\qquad z_2=-1+i,\qquad z_3=-\sqrt{3}+i,\qquad z_4=2e^{i\pi/3}.
\]

\textbf{Exercice 23 [CALC $\star\star\star$]}\\
Écrire sous forme trigonométrique puis exponentielle :
\[
z=-1+i\sqrt{3}.
\]

\textbf{Exercice 24 [CALC $\star\star\star$]}\\
Montrer que pour tous réels \(a,b\),
\[
e^{i(a+b)}=e^{ia}e^{ib}.
\]

\textbf{Exercice 25 [APP $\star\star$]}\\
Retrouver les formules d’Euler :
\[
\cos \theta=\frac{e^{i\theta}+e^{-i\theta}}{2},
\qquad
\sin \theta=\frac{e^{i\theta}-e^{-i\theta}}{2i}.
\]

\textbf{Exercice 26 [CALC $\star\star\star$]}\\
Appliquer la formule de Moivre pour calculer :
\begin{enumerate}[label=\alph*)]
    \item \((\cos\theta+i\sin\theta)^4\),
    \item \((1+i)^5\).
\end{enumerate}

\section*{VI. Équations, racines carrées et racines \(n\)-ièmes}
\addcontentsline{toc}{section}{VI. Équations, racines carrées et racines \(n\)-ièmes}

\textbf{Exercice 27 [APP $\star\star$]}\\
Déterminer les racines carrées complexes de :
\[
-4,\qquad -9,\qquad -16.
\]

\textbf{Exercice 28 [CALC $\star\star\star$]}\\
Résoudre dans \(\C\) :
\[
z^2+1=0,\qquad z^2-2z+5=0.
\]

\textbf{Exercice 29 [CALC $\star\star\star$]}\\
Résoudre dans \(\C\) :
\[
z^2-(1+4i)z+(3i-3)=0.
\]

\textbf{Exercice 30 [PREP $\star\star\star$]}\\
Déterminer les racines cubiques de l’unité.

\textbf{Exercice 31 [PREP $\star\star\star\star$]}\\
Déterminer les racines quatrièmes de \(-16\).

\textbf{Exercice 32 [PREP $\star\star\star\star$]}\\
Soit \(Z=\rho e^{i\theta}\) avec \(\rho>0\). Déterminer les solutions de
\[
z^n=Z.
\]

\section*{VII. Transformations trigonométriques}
\addcontentsline{toc}{section}{VII. Transformations trigonométriques}

\textbf{Exercice 33 [APP $\star\star\star$]}\\
Mettre sous la forme
\[
R\cos(t-\varphi)
\]
l’expression
\[
2\cos t+\sqrt{3}\sin t.
\]

\textbf{Exercice 34 [PREP $\star\star\star\star$]}\\
Mettre sous la forme
\[
R\cos(t-\varphi)
\]
puis sous la forme
\[
R\sin(t+\psi)
\]
l’expression
\[
-\cos t+2\sin t.
\]

\section*{VIII. Exercices de synthèse}
\addcontentsline{toc}{section}{VIII. Exercices de synthèse}

\textbf{Exercice 35 [PREP $\star\star\star$]}\\
Montrer que si \(|z|=1\), alors
\[
\overline z=\frac{1}{z}.
\]

\textbf{Exercice 36 [PREP $\star\star\star\star$]}\\
Soit \(z\in\C\setminus\{0\}\). Montrer que
\[
\frac{z}{\overline z}
\]
est de module \(1\).

\textbf{Exercice 37 [PREP $\star\star\star\star$]}\\
Déterminer tous les complexes \(z\) tels que
\[
z+\overline z=4
\qquad \text{et} \qquad
|z|=3.
\]

\textbf{Exercice 38 [PREP $\star\star\star\star\star$]}\\
Déterminer tous les complexes \(z\) tels que
\[
z^4=1.
\]
Les placer sur le cercle trigonométrique.

\textbf{Exercice 39 [PREP $\star\star\star\star\star$]}\\
Résoudre dans \(\C\) :
\[
z^6=-64.
\]

\textbf{Exercice 40 [PREP $\star\star\star\star\star$]}\\
Montrer que les racines \(n\)-ièmes de l’unité forment les sommets d’un polygone régulier à \(n\) côtés inscrit dans le cercle unité.

\end{document}